\section{Discussion and Conclusion}
\label{sec:conclusion}

\paragraph{What helped.}
Exemplar-free CIL methods keep class statistics only to re-train their
classifier. Using the statistics a second time, as a calibrated generative
score added to the classifier, improves RSIAT without training or new
memory. Final accuracy rises on all four benchmarks, and average accuracy on
three of them. The gain comes from correcting cross-task errors. Its size is
bounded by the representation: about $3.4$ final-accuracy points of the gap
to joint training on ImageNet-R lie in how the classes of different tasks
overlap, which no inference-time rule can change.

\paragraph{What did not.}
\begin{itemize}
  \item \emph{Transfer.} The pre-registered transfer test failed because
    CUB-200 gains no average accuracy; its gain appears only in the last
    stages.
  \item \emph{Images per class.} Our expectation that the effect grows with
    fewer images per class was wrong.
  \item \emph{Quantum components.} None of the five quantum components
    outperformed its classical counterpart. For the measurement-based
    components this follows from theory: a circuit placed after the encoding
    cannot change the PGM, and on second moments the PGM is the quantum form
    of least-squares discrimination. For the kernel-based components the
    mechanisms we measured were concentration of local reduced states and
    degenerate pseudo-data.
\end{itemize}

\paragraph{Limitations.}
\begin{itemize}
  \item We study one baseline, four of RSIAT's six benchmarks and one split
    per benchmark, with three seeds each.
  \item The fusion weights were fixed on ImageNet-R. Their use with other
    backbones or other statistics-based methods is untested.
  \item Our quantum circuits are small and simulated. The negative results
    concern these placements and encodings, not quantum models in general.
\end{itemize}

\paragraph{Outlook.}
Since the representation is the bottleneck, the next step is to train the
adapter so that the stored statistics stay discriminative, for example with
a ridge- or measurement-based loss during the incremental tasks. The
classical twin of any quantum component should be fixed as the control from
the start, as done here.

\paragraph{Reproducibility.}
Configurations, logs, analysis scripts and the pre-registered specifications
of every round are in the project repository under \texttt{exps/}. Every
number in this paper is generated from these files.
